\subsection{Coverings of topological groups}

PROPOSITION 4. --- Let (X, e) be a Hopf space (IV, p.~373, definition 2) and let m: X × X → X be its composition law. Suppose the space X is connected and locally path-connected. Let X' be a connected covering of X; denote its projection by p and let e' be a point of the fiber \(X'_{e}\).

\begin{enumerate}
\def\labelenumi{\alph{enumi})}
\item
  The covering X' is Galois and its automorphism group is commutative.
\item
  There exists a unique continuous composition law \(m' \colon X' \times X' \to X'\) such that one has \(p \circ m' = m \circ (p, p)\) and \(m'(e', e') = e'\) . Endowed with this composition law, the pointed space \((X', e')\) is a Hopf space.
\item
  Endowed with the composition law induced by \(m'\), the fiber \(X'_e\) is a group with identity element \(e'\). The map from \(\pi_1(X, x)\) into \(X'_e\) given by \(\gamma \mapsto e' \cdot \gamma\) is a surjective group homomorphism with kernel \(p_*(\pi_1(X', e'))\) . The map \(g \mapsto g(e')\) is an isomorphism of the group \(\text{Aut}_X(X')\) onto the group \(X'_e\) .
\item
  The group \(\pi_{1}(X,e)\) is commutative (prop. 3). Since the covering \(X'\) of X is connected, it is then a Galois covering and the group \(\operatorname{Aut}_{\mathrm{X}}(\mathrm{X}')\) is isomorphic to the quotient group \(\pi_{1}(\mathrm{X},e)/p_{*}(\pi_{1}(\mathrm{X}',e'))\) (III, p.~312, corollary 4 of proposition 2); this group is commutative.
\item
  Denote by \(q\colon X^{\prime}\times X^{\prime}\to X\) the map \(m\circ (p,p)\) . The map \(m'\) required is a continuous lifting of \(q\) to \(X^{\prime}\) such that \(m'(e',e') = e'\) . The space \(X^{\prime}\times X^{\prime}\) is locally path-connected; it therefore suffices, in order to prove the existence of such a continuous lifting, to verify that \(q_{*}(\pi_{1}(X^{\prime}\times X^{\prime},(e^{\prime},e^{\prime})))\) is contained in \(p_{*}(\pi_{1}(X^{\prime},e^{\prime}))\) (III, p.~308, prop. 1). By prop. 3, the map \(\pi_1(m,(e,e))\) is identified with the composition law of the group \(\pi_1(X,e)\) when one identifies \(\pi_1(X\times X,(e,e))\) with \(\pi_1(X,e)\times \pi_1(X,e)\) . The group \(q_{*}(\pi_{1}(X^{\prime}\times X^{\prime},(e^{\prime},e^{\prime})))\) is therefore equal to the group \(p_{*}(\pi_{1}(X^{\prime},e^{\prime}))\) , whence the existence of \(m'\) . Since \(X^{\prime}\times X^{\prime}\) is connected, the uniqueness of \(m'\) follows from I, p.~34, corollary 1 of prop. 11.
\end{enumerate}

Let us prove that the map \(m'\) endows the pointed space \((\mathrm{X}', e')\) with a Hopf space structure. Let \(m_1 \colon X \to X\) be the map \(g \mapsto m(g, e)\) and let \(\sigma_1 \colon X \times I \to X\) be a pointed homotopy at x connecting \(m_1\) to \(Id_X\) . Set \(\tau = \sigma_1 \circ (p, \mathrm{Id_I}) \colon X' \times I \to X\) . The map \(m'_1 \colon X' \to X'\) defined by \(h \mapsto m'(h, e')\) lifts the map \(\tau(\cdot, 0)\) ; let \(\tau' \colon X' \times I \to X'\) be the lifting of \(\tau\) whose initial map is \(m'_1\) (III, p.~301, prop. 2). The map \(t \mapsto \tau'(e', t)\) is a lift to \(X'\) of the constant map \(t \mapsto e\) ; since \(\tau'(e', 0) = e'\) , we therefore have \(\tau'(e', t) = e'\) for all \(t \in I\) . The map \(\tau'(\cdot, 1)\) is then a lifting to \(X'\) of the map p which maps \(e'\) onto \(e'\) . Since \(X'\) is connected, \(Id_{X'}\) is the unique X-morphism from \(X'\) to itself that fixes the point \(e'\) ; one therefore has \(\tau'(\cdot,1)=\mathrm{Id}_{\mathrm{X}'}\) . This shows that \(\tau'\) is a pointed homotopy at \(e'\) which connects the map \(m_{1}'\) to the map \(Id_{X'}\) . Likewise, there exists a pointed homotopy at \(e'\) connecting the map \(m_{2}'\colon h\mapsto m'(e',h)\) to the map \(Id_{X'}\) . This proves that the pointed space ( \(X',e'\) ), equipped with the composition law \(m'\) , is a Hopf space.

\begin{enumerate}
\def\labelenumi{\alph{enumi})}
\setcounter{enumi}{2}
\tightlist
\item
  Since \(X'\) is a connected covering of X, the orbit map of \(e'\) induced by the action of \(\pi_{1}(\mathrm{X},e)\) on \(X'_{e}\) is surjective and induces a bijection from \(\pi_{1}(\mathrm{X},e)/p_{*}(\pi_{1}(\mathrm{X}',e'))\) onto \(X'_{e}\) (III, p.~305, Theorem 1).
\end{enumerate}

Let \(c\) and \(d\) be loops of X at \(e\) , and let \(\gamma\) and \(\delta \in \pi_1(\mathrm{X},e)\) be their classes. Let \(c'\) and \(d'\) be the paths with initial point \(e'\) on \(\mathrm{X}'\) which lift \(c\) and \(d\) . We have \(e' \cdot \gamma = c'(1)\) and \(e' \cdot \delta = d'(1)\) (III, p.~304). By Proposition 3 of IV, p.~374, the loop \(m \circ (c,d)\) is strictly homotopic to the loop \(c*d\) ; one therefore has \(e' \cdot (\gamma\delta) = e' \cdot (m \circ (c,d))\) . Now, the path \(m' \circ (c',d')\) is a lifting with origin \(e'\) of the path \(m \circ (c,d)\) ; one therefore has

\[
e ^ {\prime} \cdot (\gamma \delta) = m ^ {\prime} (c ^ {\prime} (1), d ^ {\prime} (1)) = m ^ {\prime} (e ^ {\prime} \cdot \gamma , e ^ {\prime} \cdot \delta).
\]

The map \(\gamma \mapsto e' \cdot \gamma\) is therefore a homomorphism from \(\pi_1(\mathrm{X}, e)\) into the set \(\mathrm{X}'_e\) endowed with the composition law induced by \(m'\) . Consequently, \(\mathrm{X}'_e\) is a group under the composition law induced by \(m'\) and the orbit map of \(e'\) is an isomorphism from the quotient group \(\pi_1(\mathrm{X}, e)/p_*(\pi_1(\mathrm{X}', e'))\) onto \(\mathrm{X}'_e\) .

The last part of assertion c) then follows from Corollary 3 of III, p.~311.

PROPOSITION 5. --- We retain the notation and hypotheses of Proposition 4.

\begin{enumerate}
\def\labelenumi{\alph{enumi})}
\item
  If m is an associative (respectively commutative) law of composition, then so is \(m'\) .
\item
  If e is a right identity element (resp. left identity element) for the operation m, then e' is a right identity element (resp. left identity element) for the operation m'.
\item
  If X is a topological group, m its composition law, and e its identity element, the composition law m' endows X' with a group structure compatible with the topology of X', with e' as its identity element. The map p: X' → X is a group homomorphism whose kernel is discrete and contained in the center of X'.
\item
  Suppose that the law m is associative. Then the maps from \(X' \times X' \times X'\) into X that send \((h_{1}, h_{2}, h_{3})\) to \(m(p(h_{1}),m(p(h_{2}),p(h_{3})))\) and \(m(m(p(h_{1}),p(h_{2})),p(h_{3}))\), respectively, are equal. The maps \((h_{1},h_{2},h_{3})\mapsto m'(h_{1},m'(h_{2},h_{3}))\) and \((h_{1},h_{2},h_{3})\mapsto m'(m'(h_{1},h_{2}),h_{3})\) from \(X'\times X'\times X'\) into \(X'\) are continuous lifts that coincide at the point \((e',e',e')\) . Since \(X'\times X'\times X'\) is connected, they are equal (I, p.~34, cor. 1 of prop. 11), which shows that the law \(m'\) is associative.
\end{enumerate}

Suppose now that the law \(m\) is commutative; the maps \((h_1, h_2) \mapsto m'(h_1, h_2)\) and \((h_1, h_2) \mapsto m'(h_2, h_1)\) are continuous lifts to \(X'\) of the map \((h_1, h_2) \mapsto m(p(h_1), p(h_2))\) which coincide at the point \((e', e')\) . Since \(X' \times X'\) is connected, they are equal (loc. cit.) and the law \(m'\) is commutative.

If \(e\) is a right (resp. left) identity element for the law \(m\) , the map \(h \mapsto m'(h, e')\) (resp. \(h \mapsto m'(e', h)\) ) from \(X'\) into \(X'\) is an X-morphism of coverings that coincides with \(\text{Id}_{X'}\) at the point \(e'\) , hence at every point of \(X'\) , since \(X'\) is connected (loc. cit.), whence \(b\) ).

Let us finally prove c). Suppose that X is a topological group. By the preceding, the law \(m'\) is associative and \(e'\) is a neutral element. Let us denote by \(i: X \to X\) the map \(g \mapsto g^{-1}\) . It is continuous (TG, III, p.~1) and the homomorphism \(\pi_{1}(i,e): \pi_{1}(X,e) \to \pi_{1}(X,e)\) is none other than the map \(\gamma \mapsto \gamma^{-1}\) (IV, p.~374, prop. 3). Consequently, the subgroup \((i \circ p)_{*}(\pi_{1}(X',e'))\) is equal to \(p_{*}(\pi_{1}(X',e'))\) . By prop. 1 of III, p.~308, there therefore exists a continuous map \(i': X' \to X'\) such that \(p \circ i' = i \circ p\) and \(i'(e') = e'\) . The maps \(h \mapsto m'(h,i'(h))\) and \(h \mapsto m'(i'(h),h)\) from \(X'\) into \(X'\) are liftings to \(X'\) of the constant map with image e from \(X'\) into X. They are therefore constant and their image is \(e' = m'(e',e')\) . Thus, every element h of \(X'\) is invertible, with inverse \(i'(h)\) , which shows that \(X'\) , equipped with the composition law \(m'\) , is a group. By construction of the law \(m'\) , the map p: \(X' \to X\) is a group homomorphism. Since the maps \(m'\) and \(i'\) are continuous, the group structure of \(X'\) is compatible with its topology (TG, III, p.~1). The fiber \(p^{-1}(e)\) is a discrete subgroup of \(X'\) which is contained in the center of \(X'\) (I, p.~100, cor. 3) and X is isomorphic to the quotient topological group \(X'/p^{-1}(e)\) .

COROLLARY. --- Let G be a connected and locally path-connected topological group. Let G' be a connected covering of G, let p be its projection, and let e' be an element of the fiber N of the identity element e of G. Endow G' with the unique continuous composition law m': G' × G' → G' such that \(p \circ m' = m \circ (p, p)\) and \(m'(e', e') = e'\) . If i: N → G' denotes the canonical injection, then (G', p, i) is a central extension of G by N (A, I, p.~63).

